\begin{abstract}

%Database practitioners widely employ write-ahead logging.
%Database 

   Applications widely employ file writes to store data.	
   %Using file writes to store data is widely used by applications, such as databases.
    Researchers 
   have considered optimization tactics in both hardware and software to accelerate file writes.
   For example, many applications preallocate files in advance to minimize the cost of on-demand space allocation and journaling,
   while modern solid state drives (SSDs) are equipped with a supercapacitor for the power loss protection (PLP).
   % optimizations 
   %have optimized in hardware and software levels to accelerate file writes
   However, our study shows that, in spite of preallocation-enabled stable file structure and PLP-protected low-latency SSD,
   % these joint optimizations, 
   file writes still severely suffer from the heavy {\em runtime software tax} imposed
   by the operating system (OS), particularly when the file system maps   in-file offsets to
      device block numbers.
      
   In this paper, we develop \odes that provides a new shortcut  for file writes.
   \odes is a fully {\em all-in-software} design.
   It makes use of the stable structure enabled by preallocating files and 
   reorganizes the mappings of offsets-to-block numbers in memory 
   for quick reference. It then leverages 
   the eBPF to intercept file write or fdatasync and redirects these calls 
   %leverages the eBPF to  
   %an in-memory structure holding offset-to-block numbers for quick reference. Then it 
   %helps to intercept file write or fdatasync system calls
   to deliver data directly
   %go through ioctl interfaces  
   %take a direct path built by ioctl interfaces
   from application to device by ioctl interfaces, 
   with most of the kernel layers bypassed. For non-preallocated files, \odes utilizes the {\em move-extent} feature of Ext4
   and io\_uring, resembling a preallocation-like effect to take \odes's shortcut path.
   We use \odes to substitute the path of logging writes for two typical databases, i.e., OceanBase and MySQL.
   Evaluation results show that \odes significantly improves the performances of both.
   For example, when serving a write-intensive OLTP workload,
   \odes boosts the throughput of OceanBase by 2.1$\times$ compared to using Linux's stock storage stack.
   \odes also outperforms a state-of-the-art design that  relieves the software tax with customized hardware.
   These verify the efficiency and usability of \odes.
   
\end{abstract}
      
